\chapter{Long COVID}

\section*{Definition and Overview}
Long COVID, also known as post-COVID-19 condition or post-acute sequelae of COVID-19 (PASC), is a complex, multi-system syndrome characterized by persistent symptoms, or the development of new symptoms, that continue for 12 weeks or more after an acute infection with the severe acute respiratory syndrome coronavirus 2 (SARS-CoV-2). These symptoms fluctuate in intensity, can affect almost any organ system, and cannot be explained by an alternative medical diagnosis. Long COVID can cause significant functional impairment, severely impacting a person's ability to work, study, and perform daily activities, and requires a comprehensive, multidisciplinary clinical approach.

\section*{Epidemiology}
Long COVID has emerged as a major global health challenge. While estimates vary depending on the diagnostic criteria used, studies suggest that approximately 5\% to 20\% of individuals who contract COVID-19 will experience symptoms that persist beyond 12 weeks. In Australia, this translates to tens of thousands of individuals living with the condition. Long COVID can affect people of any age, including children, but it is most frequently diagnosed in adults aged between 30 and 60. Risk factors include female sex, severe acute COVID-19 illness (especially those requiring hospitalisation or intensive care), pre-existing chronic health conditions (such as asthma, diabetes, or autoimmune diseases), and high pre-infection body mass index. However, it can and frequently does occur in young, previously healthy individuals who experienced only mild or asymptomatic acute infections.

\section*{Aetiology and Pathophysiology}
The exact biological mechanisms underlying Long COVID are still being actively researched, but it is widely accepted to be a multi-factorial process initiated by acute SARS-CoV-2 infection. The leading pathological theories include:
\begin{itemize}
    \item \textbf{Viral persistence:} the presence of viral reservoirs, viral RNA, or persistent spike protein in various body tissues (such as the gastrointestinal tract, brain, or lymph nodes), which drives chronic, low-grade immune activation.
    \item \textbf{Autoimmunity:} the development of autoantibodies targeting the body's own tissues (such as G-protein coupled receptors or endothelial cells) due to molecular mimicry triggered by the virus.
    \item \textbf{Microvascular dysfunction and thrombosis:} persistent endothelial inflammation (endotheliitis) and the formation of microscopic, fibrin-rich blood clots (microclots) that block capillaries, impairing oxygen and nutrient delivery to tissues and causing hypoxia.
    \item \textbf{Dysautonomia:} damage to or dysfunction of the autonomic nervous system, leading to impaired regulation of heart rate, blood pressure, digestion, and temperature, frequently manifesting as postural orthostatic tachycardia syndrome (POTS).
    \item \textbf{Mitochondrial dysfunction:} impairment of cellular energy production within mitochondria, leading to profound systemic energy depletion and muscle fatigue.
    \item \textbf{Neuroinflammation:} microglial activation and inflammation within the central nervous system, leading to cognitive dysfunction and altered sensory processing.
\end{itemize}

\section*{Clinical Features}
Long COVID is characterised by a wide range of symptoms that frequently fluctuate, relapse, and worsen following physical or mental exertion. The clinical features include:
\begin{itemize}
    \item \textbf{Post-exertional malaise (PEM):} the hallmark symptom, characterised by a severe, disproportionate worsening of symptoms (such as fatigue, brain fog, or pain) 12 to 48 hours after minor physical, cognitive, or emotional effort.
    \item \textbf{Profound fatigue:} persistent, debilitating exhaustion that is not relieved by rest and is distinct from ordinary tiredness.
    \item \textbf{Cognitive dysfunction ("brain fog"):} difficulties with concentration, memory lapses, slow processing speed, word-finding difficulties, and mental fatigue.
    \item \textbf{Cardiorespiratory symptoms:} shortness of breath (dyspnoea), chronic dry cough, chest tightness, chest pain, and palpitations (a racing or pounding heart).
    \item \textbf{Orthostatic intolerance:} dizziness, lightheadedness, or feeling faint when standing up, often accompanied by a rapid heart rate (tachycardia).
    \item \textbf{Neurological and sensory symptoms:} persistent loss or alteration of taste and smell (dysgeusia and anosmia), sleep disturbances (insomnia or unrefreshing sleep), headaches, and neuropathic pain (burning, tingling, or numbness).
    \item \textbf{Gastrointestinal features:} abdominal pain, bloating, nausea, diarrhoea, and new food intolerances.
    \item \textbf{Musculoskeletal pain:} widespread muscle aches (myalgia) and joint pain (arthralgia).
\end{itemize}

\section*{Differential Diagnosis}
Because Long COVID lacks a single diagnostic test, clinicians must rule out other conditions that present with similar chronic symptoms:
\begin{itemize}
    \item \textbf{Myalgic encephalomyelitis/chronic fatigue syndrome (ME/CFS):} a chronic disease that shares almost identical clinical features, including PEM and profound fatigue, and may be triggered by other viral pathogens.
    \item \textbf{Post-intensive care syndrome (PICS):} occurs in patients who survived severe hospitalisation and ventilation, presenting with muscle weakness, cognitive impairment, and psychological distress.
    \item \textbf{Thyroid disorders:} hypothyroidism or hyperthyroidism, which can cause profound fatigue, weight changes, temperature sensitivity, and palpitations.
    \item \textbf{Cardiovascular conditions:} post-viral myocarditis, pericarditis, heart failure, or coronary artery disease, which must be excluded in patients presenting with chest pain and shortness of breath.
    \item \textbf{Primary sleep disorders:} obstructive sleep apnoea or restless legs syndrome, which cause chronic daytime fatigue and cognitive impairment.
    \item \textbf{Autoimmune diseases:} conditions like systemic lupus erythematosus (SLE) or rheumatoid arthritis, which cause chronic fatigue, joint pain, and systemic inflammation.
\end{itemize}

\section*{When to Seek Medical Care}
Individuals who experience persistent or worsening symptoms beyond 4 to 12 weeks following an acute COVID-19 infection should consult a general practitioner for assessment and support.

\textbf{Call 000 (or local emergency services) for:}
\begin{itemize}
    \item Sudden, severe, or crushing chest pain or pressure
    \item Severe difficulty breathing, gasping for air, or sudden oxygen desaturation
    \item Sudden weakness, numbness, or drooping on one side of the face or body
    \item Sudden confusion, difficulty speaking, or difficulty understanding speech
    \item Fainting, loss of consciousness, or seizures
\end{itemize}

\section*{Diagnosis and Investigations}
The diagnosis of Long COVID is clinical and based on the exclusion of other plausible causes for the patient's symptoms. The diagnostic pathway includes:
\begin{itemize}
    \item \textbf{Comprehensive clinical history:} detailing the timing of the acute COVID-19 infection, the onset of persistent symptoms, the presence of post-exertional malaise (PEM), and the impact on daily function.
    \item \textbf{Physical examination:} including lying and standing blood pressure and heart rate measurements (active stand test or NASA lean test) to screen for orthostatic intolerance.
    \item \textbf{Screening blood tests:} full blood count, kidney and liver function, thyroid stimulating hormone (TSH), HbA1c, vitamins (D and B12), iron studies, and inflammatory markers (C-reactive protein) to rule out other causes of fatigue.
    \item \textbf{Cardiovascular investigations:} electrocardiogram (ECG), echocardiogram, or a 24-hour Holter monitor to investigate chest pain, palpitations, or suspected POTS.
    \item \textbf{Respiratory investigations:} chest X-ray or high-resolution CT scan of the chest, and spirometry (lung function tests) to evaluate persistent dyspnoea.
    \item \textbf{Specialist referral:} referral to a dedicated multidisciplinary Long COVID clinic, cardiologist, respiratory physician, or neurologist if specific organ pathology is suspected.
\end{itemize}

\section*{Management}
Management is individualised and focuses on symptom relief, functional adaptation, and supporting recovery:
\begin{itemize}
    \item \textbf{Pacing and energy management:} the primary management strategy for post-exertional malaise. Patients are taught to identify their "energy envelope" and balance activity with rest, avoiding the "push-crash" cycle. Graded exercise therapy (GET) is contraindicated if PEM is present.
    \item \textbf{Autonomic management:} for patients with POTS or orthostatic intolerance, management includes high fluid intake, increased dietary salt (unless contraindicated), compression garments, and medications such as beta-blockers or ivabradine.
    \item \textbf{Pharmacological symptom relief:} low-dose naltrexone (LDN) for neuroinflammation and pain, antihistamines for suspected mast cell activation, and simple analgesics or nerve pain medications.
    \item \textbf{Sleep management:} establishing strict sleep hygiene, managing pain, and using melatonin or low-dose sedating medications if necessary.
    \item \textbf{Psychological support:} supportive counselling, cognitive behavioural therapy (CBT) adapted for chronic illness coping (not as a cure), and peer support groups to help manage the psychological impact of chronic illness.
    \item \textbf{Dietary interventions:} adopting a balanced, anti-inflammatory diet (such as the Mediterranean diet) or a low-histamine diet in patients with suspected mast cell issues.
\end{itemize}

\section*{Complications}
Long COVID can lead to substantial physical, emotional, and socioeconomic complications:
\begin{itemize}
    \item Severe loss of physical function, leading to disability and dependence on others for daily care
    \item Loss of employment, financial strain, and interrupted education
    \item Social isolation and secondary mental health conditions, including severe anxiety and clinical depression
    \item Progression to long-term chronic fatigue syndromes (such as ME/CFS) or permanent autonomic dysfunction (POTS)
    \item Increased risk of cardiovascular events, including venous thromboembolism, stroke, and myocardial infarction post-infection
\end{itemize}

\section*{Prognosis and Outcomes}
The prognosis for Long COVID is highly variable. Many individuals experience a slow but steady improvement in their symptoms over a period of 6 to 12 months. However, a significant subgroup of patients experience prolonged, debilitating symptoms lasting for years, with periods of remission followed by relapse. Early diagnosis, validation of symptoms, active pacing, and avoiding physical overexertion during recovery are associated with better outcomes. Vaccination against COVID-19 prior to infection significantly reduces the risk of developing Long COVID, and receiving vaccines post-infection can occasionally lead to an improvement in symptoms.

\section*{Prevention and Self-Care}
Prevention and self-care focus on reducing the risk of contracting the virus and managing early recovery:
\begin{itemize}
    \item Stay up to date with recommended COVID-19 vaccinations and booster doses, which significantly reduce the risk and severity of Long COVID
    \item Rest thoroughly during the acute phase of any COVID-19 infection, and avoid returning to intense physical exercise, work, or cognitive activities too quickly
    \item Keep a detailed symptom and activity diary to help identify patterns, triggers, and the onset of post-exertional malaise (PEM)
    \item Practice the "three Ps" of pacing: planning activities, prioritising tasks, and pacing yourself throughout the day
    \item Maintain good hydration and nutrition, and implement a regular, relaxing bedtime routine to support sleep quality
    \item Seek medical validation and support early, rather than pushing through symptoms, to prevent clinical deterioration
\end{itemize}

\section*{Sources and Further Reading}
The following organisations provide reliable information and guidelines on Long COVID:
\begin{itemize}
    \item Healthdirect Australia. \textit{Long COVID: symptoms, recovery, and support.} https://www.healthdirect.gov.au/
    \item Royal Australian College of General Practitioners (RACGP). \textit{Clinical guidelines for post-COVID-19 conditions.} https://www.racgp.org.au/
    \item World Health Organization. \textit{Post-COVID-19 condition definition and resources.} https://www.who.int/
    \item Centers for Disease Control and Prevention. \textit{Long COVID or post-COVID conditions.} https://www.cdc.gov/
    \item National Institute for Health and Care Excellence (NICE). \textit{COVID-19 rapid guideline: managing the long-term effects.} https://www.nice.org.uk/
\end{itemize}
